% ═════════════════════════════════════════════════════════════
% Introduction
\label{sec:intro}
% ═════════════════════════════════════════════════════════════

\noindent This project addresses the numerical solution of a regularized linear least squares problem. 
Given a data matrix $X \in \R^{m \times n}$ from the ML-CUP dataset, a random response vector 
$y \in \R^n$, and a regularization parameter $\lambda > 0$, the goal is to find the vector 
$w \in \R^m$ that minimizes

\begin{equation}
\label{eq:problem}
    \min_w \norm{\Xhat w - \yhat},
\end{equation}

\noindent where the augmented matrix and vector are defined as

\begin{equation}
\label{eq:augmented}
    \Xhat = \begin{bmatrix} X^T \\ \lambda I_m \end{bmatrix} \in \R^{(n+m) \times m}, 
    \qquad
    \yhat = \begin{bmatrix} y \\ 0 \end{bmatrix} \in \R^{n+m}.
\end{equation}

\noindent Expanding the squared norm, the problem is equivalent to minimizing

\begin{equation}
\label{eq:expanded}
    f(w) = \frac{1}{2}\norm{X^T w - y}^2 + \frac{1}{2}\lambda^2 \norm{w}^2,
\end{equation}

\noindent which is a standard instance of Tikhonov regularization (ridge regression). 
The regularization term $\lambda^2\norm{w}^2$ penalizes large coefficient vectors, 
improving the numerical stability of the solution and controlling overfitting.

\noindent Two computational strategies are explored:

\begin{itemize}[nosep]
    \item \textbf{(A1)} A limited-memory quasi-Newton method (L-BFGS), which iteratively 
    approximates the curvature of the objective function using only gradient information 
    and a small number of stored vector pairs.
    \item \textbf{(A2)} A direct method based on thin QR factorization with Householder 
    reflectors, in the compact variant where the orthogonal matrix~$Q$ is not formed 
    explicitly, but applied implicitly through stored Householder vectors.
\end{itemize}

\noindent Both algorithms are implemented from scratch, without relying on any off-the-shelf 
numerical solvers.